\chapter{Formulácia úlohy}\label{ch:task-definition}

Úlohou tejto bakalárskej práce je vytvorenie dvoch funkčne ekvivalentných verzií webovej hry Dáma a následné porovnanie ich vlastností z hľadiska implementácie, prevádzky a údržby. Hlavným cieľom je prakticky demonštrovať dopady voľby architektúry na vývojový proces, nasadzovanie a dlhodobú udržateľnosť systému a poskytnúť objektívne závery pre rozhodovanie o vhodnosti architektonického prístupu pri podobných projektoch.

\section{Konkrétne požiadavky úlohy}

\subsection{Monolitická verzia}

Implementovať aplikáciu ako jeden nasaditeľný artefakt, v ktorom sú sústredené všetky komponenty – frontend, backend a perzistencia dát. Tento prístup umožňuje jednoduchšie nasadenie, rýchlu internú komunikáciu a minimálne infraštruktúrne nároky. Výzvou je udržať kód prehľadný a modulárny napriek tomu, že všetky funkcionality sú integrované v jednom celku. Aplikácia musí podporovať správu hráčov, vykonávanie ťahov, validáciu pravidiel, ukončenie partie a perzistentné ukladanie výsledkov a komentárov.

\subsection{Mikroservisná verzia}

Implementovať architektúru pozostávajúcu z viacerých nezávislých služieb, kde každá rieši konkrétnu doménu (hráči, hra, skóre, komentáre, hodnotenia). Služby komunikujú cez \glslink{restapi}{REST API} a prístup k nim je centralizovaný cez \glslink{gateway}{API Gateway}. Tento prístup zabezpečuje nezávislé nasadzovanie, lepšiu odolnosť voči chybám a možnosť škálovať jednotlivé časti systému, pričom zvyšuje komplexitu konfigurácie, orchestrácie a monitorovania.

\subsection{Funkčná ekvivalencia}

Obe riešenia musia poskytovať totožnú používateľskú skúsenosť. Frontend je v oboch prípadoch implementovaný ako \glslink{react}{React} SPA, ktorý zobrazuje hraciu dosku, spracúva ťahy, komunikuje s backendom a zobrazuje leaderboard, profil hráčov či komentáre.

\subsection{Dátový model a API kontrakty}

\glslink{monolith}{Monolit} používa jednu databázovú schému. \glslink{microservice}{Mikroservisy} využívajú spoločnú \glslink{postgresql}{PostgreSQL} databázu, avšak každá služba pracuje s vlastnými tabuľkami v rámci oddelených domén. \glslink{api}{API} kontrakty definujú presnú podobu komunikácie, formáty JSON odpovedí, chybové stavy a validačné pravidlá.

\subsection{Spustenie a testovanie}

Obe verzie musia byť spustiteľné cez \glslink{dockercompose}{Docker Compose}. Súčasťou úlohy je navrhnúť testovacie scenáre, pripraviť jednotkové a integračné testy a následne vykonať merania výkonnosti.

\subsection{Porovnávacia analýza}

Na základe testov a meraní bude vypracovaná porovnávacia štúdia, ktorá zhodnotí výhody a nevýhody oboch prístupov z hľadiska výkonu, modularity, prevádzky, škálovateľnosti a údržby.

Popri technických cieľoch má práca aj širšie vzdelávacie a metodické ambície. Jedným z čiastkových cieľov je ukázať, že architektonické rozhodnutia nemožno hodnotiť izolovane, ale vždy v kontexte konkrétneho projektu, tímu a infraštruktúrnych možností. Pre malé projekty môže byť \glslink{monolith}{monolit} s dobre navrhnutou vnútornou štruktúrou plne postačujúci, kým pre väčšie systémy s vyššími nárokmi na škálovateľnosť a nezávislé nasadzovanie je vhodnejšia \glslink{microservice}{mikroservisná architektúra}. Výsledky práce by preto mali pomôcť nielen pri voľbe konkrétnej technológie, ale aj pri formovaní spôsobu uvažovania o architektúre ako o súbore kompromisov, ktoré je potrebné explicitne pomenovať a zdôvodniť.